\subsection{Support Vector Machine}
\label{sec:support_vector_machine}

\begin{table*}[htbp]
\caption{Hyperparameter Search Space for Support Vector Machine Classifier.}
\label{tab:6}
\label{tab:VI}
\label{tab:svm_params}
\centering
\begin{tabular}{llc}
\toprule
Component / Hyperparameter & Evaluated Values & Count \\
\midrule
Feature Scaler & None, MinMaxScaler & 2 \\
Dimensionality Reduction (PCA) & None, 0.50, 0.75 & 3 \\
Regularization Parameter & 0.01, 0.1, 1, 10 & 4 \\
Kernel Function & 'linear', 'rbf', 'poly' & 3 \\
Polynomial Degree & 2, 3 & 2 \\
Kernel Coefficient & 'scale', 'auto' & 2 \\
\midrule
\multicolumn{2}{r}{\textbf{Total Grid Combinations}} & \textbf{288 (1,440 fits)} \\
\bottomrule
\end{tabular}
\end{table*}

Pengklasifikasi Support Vector Machine (SVM) mengonstruksikan bidang pemisah berjarak terjauh (maximum margin hyperplane) untuk memisahkan enam kelas interseksional pada ruang laten berdimensi tinggi~\cite{ref40}. Fungsi kernel memproyeksikan distribusi nonlinier fitur multi-domain ke ruang berdimensi lebih tinggi, mencakup kernel linear, Radial Basis Function (RBF), dan polinomial~\cite{ref41}. Formulasi kernel polinomial pada \eqref{eq:11} menghitung produk titik $\langle \mathbf{x}_i, \mathbf{x}_j \rangle$ antara vektor fitur $\mathbf{x}_i$ dan $\mathbf{x}_j$, di mana $\gamma$ menyatakan skala kernel, $d \in \{2, 3\}$ menentukan derajat pemetaan, dan konstanta intercept $\text{coef0} = 0.0$ mengikuti nilai bawaan scikit-learn~\cite{ref42}:
\begin{equation}
K(\mathbf{x}_i, \mathbf{x}_j) = (\gamma \langle \mathbf{x}_i, \mathbf{x}_j \rangle + \text{coef0})^d, \quad d \in \{2, 3\}
\label{eq:11}
\end{equation}
Pemetaan polinomial ini memfasilitasi penangkapan interaksi nonlinier antardimensi visual secara efisien~\cite{ref43}. Optimasi hyperparameter mengevaluasi parameter penalti $C \in \{0.01, 0.1, 1, 10\}$, ragam kernel, koefisien $\gamma$, derajat polinomial, serta opsi penskalaan dan PCA, menghasilkan 288 kombinasi grid Cartesian (108 konfigurasi unik efektif) atau 1.440 total fits pada 5-Fold StratifiedGroupKFold Cross-Validation dengan isolasi subjek sesuai Table~\ref{tab:6}.
